\subsection{LINEAR REPRESENTATIONS}

Let \(G\) be a Lie group, \(E\) a complete normable space and \(\pi\) an analytic linear representation of \(G\) on \(E\) (§ 1, no. 2). The associated morphism \(t \mapsto \langle t, \pi \rangle\) of \(\mathcal{F}^{(\infty)}(G)\) into \(\mathcal{L}(E)\) is an algebra morphism (no. 9, Proposition 33) and its restriction to \(L(G)\) is \(L(\pi)\) . Hence \(L(\pi)\) is a representation of \(L(G)\) on \(E\) (Chapter I, § 3, Definition 1).

PROPOSITION 38. Consider G as operating on E on the left by the mapping \((g,x)\mapsto \pi (g)x\) . Let \(b\in \mathrm{E}\) and \(\rho (b)\) be its orbital mapping. Let \(\mathbf{T}_b(\mathbf{E})\) be canonically identified with E. For all \(t\in \mathbf{L}(\mathbf{G})\)

\[
(\mathbf {L} (\pi) t) (b) = \langle t, \rho (b) \rangle = \rho (b) _ {*} t = t * \varepsilon_ {b}.
\]

In particular, the vector field defined by \(t\) on \(\mathbf{E}\) is the field \(b \mapsto (\mathbf{L}(\pi)t)(b)\) .

\[
\begin{array}{r l} (\mathrm{L} (\pi) t) (b) & = \langle t, g \mapsto \pi (g) b \rangle \\ & = \langle t, \mathrm{Id} _ {\mathrm{E}} \circ \rho (b) \rangle \\ & = \langle \rho (b) _ {*} t, \mathrm{Id} _ {\mathrm{E}} \rangle \quad (Diff. \& Anal. Man., R, 13.2.3) \\ & = \rho (b) _ {*} t. \end{array}
\]

Finally, \(\rho(b)_{*}t = t*\varepsilon_{b}\) (no. 3, Proposition 14 (ii)).

PROPOSITION 39. Suppose that K is of characteristic 0. Let G be a Lie group, E a finite-dimensional vector space and \(\pi\) an analytic linear representation of G on E. Let \(\mathrm{E}_1\) , \(\mathrm{E}_2\) be vector subspaces of E such that \(\mathrm{E}_2 \subset \mathrm{E}_1\) . The set \(\mathrm{G}_1\) of \(g \in \mathrm{G}\) such that \(\pi(g) x \equiv x \pmod{\mathrm{E}_2}\) for all \(x \in \mathrm{E}_1\) is a Lie subgroup of G and \(\mathrm{L}(\mathrm{G}_1)\) is the set of \(a \in \mathrm{L}(\mathrm{G})\) such that \(\mathrm{L}(\pi)a\) maps \(\mathrm{E}_1\) into \(\mathrm{E}_2\) .

This follows from Propositions 29 (no. 8) and 36 (no. 10).

COROLLARY 1. In the notation of Proposition 39, the set of \(g \in G\) such that \(\pi(g)(\mathrm{E}_1) \subset \mathrm{E}_1\) is a Lie subgroup of \(G\) and its Lie algebra is the set of \(a \in L(G)\) such that \(L(\pi)a\) maps \(E_1\) into \(E_1\) .

We apply Proposition 39 with \(E_{1} = E_{2}\) .

COROLLARY 2. Let G, E, \(\pi\) be as in Proposition 39. Let F be a subset of E. The set of \(g \in \mathbf{G}\) such that \(\pi(g)x = x\) for all \(x \in \mathbf{F}\) is a Lie subgroup of G and its Lie algebra is the set of \(a \in \mathbf{L}(\mathbf{G})\) such that \((\mathbf{L}(\pi)a)(x) = 0\) for all \(x \in \mathbf{F}\) .

We apply Proposition 39 with \(E_{2}=\{0\}\) and \(E_{1}\) the vector subspace of E generated by F.

Let \(\pi_1, \pi_2, \ldots, \pi_n\) be analytic linear representations of \(G\) . Clearly the direct sum \(\pi\) of the \(\pi_1\) (Algebra, Chapter VIII, § 13, no. 1) is an analytic linear representation of \(G\) and \(L(\pi)\) is the direct sum of \(L(\pi_1), L(\pi_2), \ldots, L(\pi_n)\) (Chapter I, § 3, no. 1).

PROPOSITION 40. Let G be a Lie group, E a complete normable space, \(\pi\) an analytic linear representation of G on E and F a closed vector subspace of E stable under \(\pi(\mathbf{G})\) . Suppose that K is of characteristic 0, or that F is a direct factor of E.

\begin{enumerate}
\def\labelenumi{(\roman{enumi})}
\item
  The subrepresentation \(\pi_1\) and the quotient representation \(\pi_2\) of \(\pi\) defined by F are analytic representations.
\item
  F is stable under L(π)(L(G)).
\item
  Let \(\rho_{1}\) and \(\rho_{2}\) be the subrepresentation and quotient representation of \(\mathbf{L}(\pi)\) defined by F. Then \(\mathbf{L}(\pi_1) = \rho_1, \mathbf{L}(\pi_2) = \rho_2\) .
\end{enumerate}

Let \(\mathrm{A}\) be the set of \(u\in \mathcal{L}(\mathrm{E})\) such that \(u(\mathrm{F})\subset \mathrm{F}\) . Then \(\mathrm{A}\) is a closed vector subspace of \(\mathcal{L}(\mathrm{E})\) and \(\pi\) takes its values in \(\mathbf{A}\) . By virtue of the hypotheses on \(\mathbf{K}\) and \(\mathbf{F}\) , the mapping \(\pi^{\prime}\colon \mathrm{G}\to \mathrm{A}\) with the same graph as \(\pi\) is analytic (Differentiable and Analytic Manifolds, R, 5.8.5). The canonical mappings \(\theta_{1}\colon \mathrm{A}\to \mathcal{L}(\mathrm{F})\) and \(\theta_{2}\colon \mathrm{A}\to \mathcal{L}(\mathrm{E / F})\) are continuous and linear and hence analytic. This proves (i). The mappings \(\mathrm{T_e}(\pi)\) and \(\mathrm{T_e}(\pi ')\) have the same graph and hence \(\mathrm{L}(\pi)(\mathrm{L}(\mathrm{G}))\subset \mathrm{A}\) , which proves (ii). We have

\[
\begin{array}{l} \mathrm{T} _ {e} (\pi_ {1}) = \mathrm{T} _ {e} (\theta_ {1} \circ \pi^ {\prime}) = \theta_ {1} \circ \mathrm{T} _ {e} (\pi^ {\prime}) = \rho_ {1} \\ \mathrm{T} _ {e} (\pi_ {2}) = \mathrm{T} _ {e} (\theta_ {2} \circ \pi^ {\prime}) = \theta_ {2} \circ \mathrm{T} _ {e} (\pi^ {\prime}) = \rho_ {2}. \end{array}
\]

PROPOSITION 41. Let G be a Lie group and \(\pi_1, \pi_2, \ldots, \pi_n\) , \(\pi\) analytic linear representation of G on complete normable spaces \(E_1, E_2, \ldots, E_n\) , E. Let

\[
\left(x _ {1}, x _ {2}, \dots , x _ {n}\right) \mapsto x _ {1} x _ {2} \dots x _ {n}
\]

be a continuous multilinear mapping of \(E_{1} \times E_{2} \times \cdots \times E_{n}\) into E. Suppose that

\[
\pi (g) \left(x _ {1} x _ {2} \dots x _ {n}\right) = \left(\pi_ {1} (g) x _ {1}\right) \left(\pi_ {2} (g) x _ {2}\right) \dots \left(\pi_ {n} (g) x _ {n}\right)
\]

for all \(g \in \mathbf{G}\) , \(x_{1} \in \mathrm{E}_{1}, \ldots, x_{n} \in \mathrm{E}_{n}\) . Then

\[
(\mathrm{L} (\pi) a) (x _ {1} x _ {2} \dots x _ {n}) = \sum_ {i = 1} ^ {n} x _ {1} x _ {2} \dots x _ {i - 1} ((\mathrm{L} (\pi_ {i}) a) x _ {i}) x _ {i + 1} \dots x _ {n}
\]

for all \(a \in \mathrm{L}(\mathrm{G})\) , \(x_1 \in \mathrm{E}_1, \ldots, x_n \in \mathrm{E}_n\) .

As an example we perform the calculation for n = 2.

\[
\begin{array}{l l} (\mathrm{L} (\pi) a) (x _ {1} x _ {2}) & = \langle a, g \mapsto \pi (g) (x _ {1} x _ {2}) \rangle \\ & = \langle a, (g \mapsto \pi_ {1} (g) x _ {1}) (g \mapsto \pi_ {2} (g) x _ {2}) \rangle \\ & = \langle a, g \mapsto \pi_ {1} (g) x _ {1} \rangle . x _ {2} + x _ {1}. \langle a, g \mapsto \pi_ {2} (g) x _ {2} \rangle \\ & = ((\mathrm{L} (\pi_ {1}) a) x _ {1}). x _ {2} + x _ {1}. ((\mathrm{L} (\pi_ {2}) a) x _ {2}) \quad \text {(Proposition 38).} \end{array}
\]

COROLLARY 1. Let G be a Lie group, \(\mathbf{E}_1,\ldots ,\mathbf{E}_{n + 1}\) complete normable spaces and \(\pi_1,\dots ,\pi_{n + 1}\) analytic linear representations of G on \(\mathbf{E}_1,\dots ,\mathbf{E}_{n + 1}\) . Let

\[
\mathrm{E} = \mathscr {L} (\mathrm{E} _ {1}, \dots , \mathrm{E} _ {n}; \mathrm{E} _ {n + 1})
\]

the complete normable space of continuous multilinear mappings of \(E_{1} \times \cdots \times E_{n}\) into \(E_{n+1}\) (General Topology, Chapter X, §3, no. 2). For all \(g \in G\) , let \(\pi(g)\) be the automorphism of E defined by

\[
(\pi (g) u) (x _ {1}, \dots , x _ {n}) = \pi_ {n + 1} (g) (u (\pi_ {1} (g) ^ {- 1} x _ {1}, \dots , \pi_ {n} (g) ^ {- 1} x _ {n})).
\]

Then \(\pi\) is an analytic linear representation of \(\mathbf{G}\) on \(\mathbf{E}\) and

\[
\begin{array}{l} ((\mathrm{L} (\pi) a) u) (x _ {1}, \ldots , x _ {n}) = - \sum_ {i = 1} ^ {n} u (x _ {1}, \ldots , x _ {i - 1}, (\mathrm{L} (\pi_ {i}) a) x _ {i}, x _ {i + 1}, \ldots , x _ {n}) \\ \qquad + (\mathrm{L} (\pi_ {n + 1}) a) (u (x _ {1}, \ldots , x _ {n})) \end{array}
\]

for all \(a \in \mathrm{L}(\mathrm{G})\) , \(u \in \mathrm{E}\) , \(x_1 \in \mathrm{E}_1, \ldots, x_n \in \mathrm{E}_n\) .

Every element \((A_{1},\ldots,A_{n+1})\) of \(\mathcal{L}(\mathrm{E}_{1})\times\cdots\times\mathcal{L}(\mathrm{E}_{n+1})\) defines a continuous endomorphism \(\theta(A_{1},\ldots,A_{n+1})\) of E by the formula

\[
(\theta (A _ {1}, \dots , A _ {n + 1}) u) (x _ {1}, \dots , x _ {n}) = A _ {n + 1} (u (A _ {1} x _ {1}, \dots , A _ {n} x _ {n})).
\]

The mapping \(\theta\) of \(\mathcal{L}(\mathrm{E}_1) \times \cdots \times \mathcal{L}(\mathrm{E}_{n+1})\) into \(\mathcal{L}(\mathrm{E})\) is continuous and multilinear. Then, for all \(g \in \mathrm{G}\) ,

\[
\pi (g) = \theta (\pi_ {1} (g ^ {- 1}), \dots , \pi_ {n} (g ^ {- 1}), \pi_ {n + 1} (g))
\]

and hence \(\pi\) is analytic. We apply Proposition 41 to the mapping

\[
\left(x _ {1}, \dots , x _ {n}, u\right) \mapsto u \left(x _ {1}, \dots , x _ {n}\right)
\]

of \(\mathbf{E}_1\times \dots \times \mathbf{E}_n\times \mathbf{E}\) into \(\mathbf{E}_{n + 1}\) . Then

\[
\pi_ {n + 1} (g) \left(u \left(x _ {1}, \dots , x _ {n}\right)\right) = (\pi (g) u) \left(\pi_ {1} (g) x _ {1}, \dots , \pi_ {n} (g) x _ {n}\right)
\]

and hence

\[
\begin{array}{l} (\mathrm{L} (\pi_ {n + 1}) a) (u (x _ {1}, \ldots , x _ {n})) = \sum_ {i = 1} ^ {n} u (x _ {1}, \ldots , (\mathrm{L} (\pi_ {i}) a) x _ {i}, \ldots , x _ {n}) \\ \qquad \qquad \qquad + ((\mathrm{L} (\pi) a) u) (x _ {1}, \ldots , x _ {n}). \end{array}
\]

When the \(E_{i}\) are finite-dimensional, the representation \(\mathrm{L}(\pi)\) of \(\mathrm{L}(\mathrm{G})\) is derived from the representations \(\mathrm{L}(\pi_{1}),\ldots,\mathrm{L}(\pi_{n+1})\) under the procedure of Chapter I, §3, Proposition 3.

COROLLARY 2. Let G be a Lie group and \(\pi\) an analytic linear representation of G on a complete normable space E. Then \(g \mapsto {}^{t}\pi(g)^{-1}\) is an analytic linear representation \(\rho\) of G on the complete normable space \(\mathcal{L}(\mathrm{E}, \mathrm{K})\dagger\) and \(\mathrm{L}(\rho)a = -{}^{t}(\mathrm{L}(\pi)a)\) for all \(a \in \mathrm{L}(\mathrm{G})\) .

This is a special case of Corollary 1.

\(\rho\) is called the contragredient representation of \(\pi\) .

When E is finite-dimensional, \(\mathbf{L}(\rho)\) is the dual representation of \(\mathbf{L}(\pi)\) in the sense of Chapter I, § 3, no. 3.

COROLLARY 3. Let G be a Lie group and \(\pi_1, \ldots, \pi_n\) analytic linear representations of G on finite-dimensional vector spaces \(E_{1},\ldots,E_{n}\) . Then the representation \(\pi_{1}\otimes\cdots\otimes\pi_{n}\) of G (Appendix) is analytic and \(\mathrm{L}(\pi_{1}\otimes\cdots\otimes\pi_{n})\) is the tensor product of \(\mathrm{L}(\pi_{1}),\ldots,\mathrm{L}(\pi_{n})\) .

The mapping \((A_{1},\ldots,A_{n})\mapsto A_{1}\otimes\cdots\otimes A_{n}\) of \(\mathcal{L}(\mathrm{E}_{1})\times\cdots\times\mathcal{L}(\mathrm{E}_{n})\) into \(\mathcal{L}(\mathrm{E}_{1}\otimes\cdots\otimes\mathrm{E}_{n})\) is multilinear, whence the fact that \(\pi\) is analytic. Consider the mapping \((x_{1},\ldots,x_{n})\mapsto x_{1}\otimes\cdots\otimes x_{n}\) of \(E_{1}\times\cdots\times E_{n}\) into

\[
\mathrm{E} _ {1} \otimes \dots \otimes \mathrm{E} _ {n}.
\]

By Proposition 41, we see that

\[
(\mathrm{L} (\pi) a) (x _ {1} \otimes \dots \otimes x _ {n}) = \sum_ {i = 1} ^ {n} x _ {1} \otimes \dots \otimes (\mathrm{L} (\pi_ {i}) a) x _ {i} \otimes \dots \otimes x _ {n}
\]

for all \(a \in \mathrm{L}(\mathrm{G})\) , \(x_{i} \in \mathrm{E}_{i}\) for \(1 \leqslant i \leqslant n\) . Hence \(\mathrm{L}(\pi)\) is the tensor product of the \(\mathrm{L}(\pi_{i})\) .

COROLLARY 4. Let G be a Lie group and \(\pi\) an analytic linear representation of G on a finite-dimensional vector space E. Then the representations \(T^{n}(\pi)\) , \(S^{n}(\pi)\) and \(\wedge^{n}(\pi)\) of G (Appendix) are analytic and

\[
\mathrm{L} \left(\mathrm{T} ^ {n} (\pi)\right) = \mathrm{T} ^ {n} (\mathrm{L} (\pi)), \quad \mathrm{L} \left(\mathrm{S} ^ {n} (\pi)\right) = \mathrm{S} ^ {n} (\mathrm{L} (\pi)), \quad \mathrm{L} \left(\wedge^ {n} (\pi)\right) = \wedge^ {n} (\mathrm{L} (\pi)).
\]

This follows from Corollary 3 and Proposition 40.

COROLLARY 5. Let A be a finite-dimensional algebra. Suppose that K is of characteristic 0. The automorphism group \(\operatorname{Aut}(\mathbf{A})\) of A is a Lie subgroup of \(\mathbf{GL}(\mathbf{A})\) and \(\mathbf{L}(\operatorname{Aut}(\mathbf{A}))\) is the Lie algebra of derivation of A.

This follows from Corollary 1 (applied to \(\mathbf{E} = \mathcal{L}(\mathbf{A},\mathbf{A};\mathbf{A})\) ) and Corollary 2 of Proposition 39 (applied to the subset of \(\mathbf{E}\) consisting only of multiplication on \(\mathbf{A}\) ).

Remark. We apply Corollary 1 with \(\mathrm{G} = \mathbf{GL}(\mathrm{F})\) (F a complete normable space), \(\pi_1 = \pi_2 = \mathrm{Id}_{G}\) and \(\pi_3\) the trivial representation of G on K. We obtain an analytic representation \(\pi\) of \(\mathbf{GL}(\mathrm{F})\) on \(\mathcal{L}(\mathrm{F},\mathrm{F};\mathrm{K})\) . We assume that F is finite-dimensional and that K is of characteristic 0. Applying Corollary 2 to Proposition 39 to \(\pi\) , we recover part of Corollary 1 to Proposition 37.

PROPOSITION 42. Let G be a Lie group, X an analytic manifold, \((g,x)\mapsto gx\) (resp. \(xg\) ) a law of analytic left (resp. right) operation of G on X and \(x_0\) a point of X which is invariant under G. For all \(g\in G\) , let \(\tau (g)\) be the automorphism \(x\mapsto gx\) (resp. \(xg\) ) of X and let \(\pi (g)\) be the automorphism of \(\mathrm{T}_{x_0}(\mathbf{X})\) tangent at \(x_0\) to \(\tau (g)\) .

\begin{enumerate}
\def\labelenumi{(\roman{enumi})}
\item
  \(\pi\) is an analytic representation of G (resp. \(\mathbf{G}^{\vee}\) ) on \(\mathrm{T}_{x_0}(\mathbf{X})\) .
\item
  For all \(a \in \mathbf{L}(\mathbf{G})\) and all \(\xi_0 \in \mathrm{T}_{x_0}(\mathbf{X}), \mathrm{L}(\pi)a\) . \(\xi_0\) can be calculated as follows: let \(\mathbf{D}_a\) be the vector field defined by \(a\) on \(\mathbf{X}\) and \(\xi\) a vector field of class \(\mathbf{C}^1\) in an open neighbourhood of \(x_0\) such that \(\xi(x_0) = \xi_0\) ; then
\end{enumerate}

\[
\mathrm{L} (\pi) a. \xi_ {0} = - [ \mathrm{D} _ {a}, \xi ] (x _ {0}).
\]

\(\tau(gg') = \tau(g)\tau(g')\) (resp. \(\tau(g')\tau(g)\) ) and hence \(\pi(gg') = \pi(g)\pi(g')\) (resp. \(\pi(g')\pi(g)\) ). On the other hand, since TX is a vector G-bundle of class \(\mathbf{C}^{\omega}\) (§ 1, no. 8, Proposition 16), \(\pi\) is analytic, whence (i).

To prove (ii), suppose that \(G\) operates on the left. There exists an open neighbourhood \(I\) of \(0\) in \(K\) and an analytic mapping \(\gamma\) of \(I\) into \(G\) such that \(\gamma(0) = e\) , \(T_0(\gamma)1 = a\) . Then \(D_a\) is the vector field on \(X\) defined by the mapping \(\phi: (\lambda, x) \mapsto \gamma(\lambda)x\) of \(I \times X\) into \(X\) (§ 2, no. 2). If \(\phi_{\lambda}\) denotes the bijection \(x \mapsto \gamma(\lambda)x\) of \(X\) into \(X\) , then

\[
\begin{array}{r l} \left[ D _ {a}, \xi \right] (x _ {0}) & = \left(\frac {d}{d \lambda} \left(T _ {\phi_ {\lambda} (x _ {0})} (\phi_ {\lambda} ^ {- 1}) \xi (\phi_ {\lambda} (x _ {0}))\right)\right) _ {\lambda = 0} \quad \text {(Diff. \& Anal. Man., R, 8.4.5)} \\ & = \left(\frac {d}{d \lambda} \left(T _ {x _ {0}} (\phi_ {\lambda} ^ {- 1}) \xi_ {0}\right)\right) _ {\lambda = 0} \\ & = \left(\frac {d}{d \lambda} \left(\pi (\gamma (\lambda)) ^ {- 1} \xi_ {0}\right)\right) _ {\lambda = 0}. \end{array}
\]

As the mappings \(\lambda \mapsto \gamma (\lambda)^{-1}\) and \(\lambda \mapsto \gamma (-\lambda)\) are tangent at 0, this is also equal to

\[
\begin{array}{r l} & {- \left(\frac {d}{d \lambda} (\pi (\gamma (\lambda)) \xi_ {0})\right) _ {\lambda = 0}} \\ {=} & {- \left(\frac {d}{d \lambda} (\pi \circ \gamma) (\lambda)\right) _ {\lambda = 0} \xi_ {0}} \\ {=} & {- \mathrm{L} (\pi) a. \xi_ {0}.} \end{array}
\]
% semantic-fragments: legacy-input-seam
\relax{}
